\documentclass[noauthor]{ximera}

\title{Week 5}
\author{Math 1193 Design Team}

\begin{document}
\begin{abstract}
    Solutions for Week 5.
\end{abstract}
\maketitle

\begin{problem}
There is a special point on the graph of a parabola called the 
\textbf{vertex}. For the parent function \(f\left(x\right)=x^2\), 
the vertex is located at the point \((0,0)\). Each transformation 
applied to the parent function moves the vertex.

\begin{enumerate}
    \item What is the vertex of the graph of \(g(x) = (x - 2)^2\)?
    \item What is the vertex of the graph of \(g(x) = x^2 - 2\)?
    \item What is the vertex of the graph of \(g(x) = 2x^2\)?
    \item What is the vertex of the graph of \(g(x) = (-x)^2\)?
\end{enumerate}

\begin{explanation}
\begin{enumerate}
    \item \((2, 0)\)
    \item \((0, -2)\)
    \item \((0, 0)\)
    \item \((0, 0)\)
\end{enumerate}
\end{explanation}
\end{problem}

\begin{problem}
\begin{enumerate}
    \item Write down the equation of a parabola whose graph has
    its vertex at \((-1, 3)\).
    \item Can you find another?
    \item How many such parabolas are there?
\end{enumerate}

\begin{explanation}
\begin{enumerate}
    \item \(f(x) = (x + 1)^2 + 3\)
    \item \(f(x) = -(x + 1)^2 + 3\) is one possibility.
    \item There are infinitely many such parabolas.
\end{enumerate}
\end{explanation}
\end{problem}

\begin{problem}
A quadratic function is one of the form \(ax^2 + bx + c\), where \(a\),
\(b\), and \(c\) are real numbers (\(b\) and \(c\) may be 0). Select all the 
following functions which are quadratic. 

For each quadratic function below, find its zeros (zeros are 
\(x\)-values where the functions output is 0). Check with your group.
Did you all use the same method? Can you understand each other's methods?

\begin{enumerate}
    \item \(g\left(x\right)=x^2-9\)
    \item \(f\left(x\right)=6+5x+x^2\)
    \item \(g\left(x\right)=0.1x^2\)
    \item \(h\left(x\right)=\left(x-1\right)\left(x-2\right)^2\)
    \item \(f\left(x\right)=2\left(x+2\right)\left(x-9\right)\)
    \item \(f\left(x\right)=\left(3x-1\right)\left(x+1\right)\)
    \item \(g\left(x\right)=-3\left(x-\frac{1}{2}\right)^2+4\)
    \item \(h\left(x\right)=x\left(x+3\right)\)
    \item \(f\left(x\right)=2\left(x-6\right)\)
    \item \(g\left(x\right)=2x^2-11x+5\)
\end{enumerate}

\begin{explanation}
\begin{enumerate}
    \item \(g\) is quadratic. It doesn't have an \(x\) term, but that's okay (we can view the \(x\) as having 
    a coefficient of \(0\)).
    \begin{align*}
        x^2 - 9 & = 0 \\
        x^2 & = 9 \\
        x & = \pm 3
    \end{align*}
    is one possible way to find the zeros. Remember that when taking an even
    root on both sides of an equation, you need to add a \(\pm\).

    \item \(f\) is quadratic. The order of addition doesn't matter, and the highest power of \(x\) is 2.
    \begin{align*}
        x^2 + 5x + 6 & = 0 \\
        (x + 3)(x + 2) & = 0 \\
    \end{align*}
    meaning \(x + 3 = 0\) or \(x + 2 = 0\), so \(x = -3, -2\). 
    
    \item \(g\) is quadratic. Again, the constant and \(x\) terms are allowed to have a coefficient of 0.
    \begin{align*}
        0.1x^2 & = 0 \\
        x^2 & = 0 \\
        x & = \pm 0
    \end{align*}
    so \(x = 0\).

    \item \(h\) is not quadratic: when you multiply the factors, you will obtain an \(x^3\) term.
    
    \item \(f\) is quadratic. Multiplying the factors reveals an \(x^2\) term.
    \begin{align*}
        2\left(x+2\right)\left(x-9\right) & = 0 \\
        (x + 2)(x - 9) & = 0 \\
    \end{align*}
    meaning \(x + 2 = 0\) or \(x - 9 = 0\), so \(x = -2, 9\).

    \item \(f\) is quadratic.
    \begin{align*}
        \left(3x-1\right)\left(x+1\right) & = 0
    \end{align*}
    so \(3x - 1 = 0\) or \(x + 1 = 0\). After rearranging, this means \(x = -1, \frac{1}{3}\).

    \item \(g\) is quadratic. Squaring the \((x - 1/2)\) will give you an \(x^2\) term.
    \begin{align*}
        -3\left(x-\frac{1}{2}\right)^2+4 & = 0 \\
        -3\left(x-\frac{1}{2}\right)^2 & = -4 \\
        \left(x-\frac{1}{2}\right)^2 & = \frac{4}{3} \\
        x - \frac{1}{2} & = \pm \sqrt{\frac{4}{3}} \\
        x & = \pm \frac{2}{\sqrt{3}} + \frac{1}{2}
    \end{align*}

    \item \(h\) is quadratic. When you distribute, you obtain an \(x^2\) term.
    \begin{align*}
        x\left(x+3\right) & = 0 \\
    \end{align*}
    so \(x = 0\) or \(x + 3 = 0\), meaning \(x = -3, 0\).

    \item \(f\) is not quadratic. When you distribute, you do not obtain an \(x^2\) term.
    \item \(g\) is quadratic. 
    \begin{align*}
        2x^2-11x+5 & = 0 \\
        (2x - 1)(x - 5) & = 0
    \end{align*}
    meaning \(2x - 1 = 0\) or \(x - 5 = 0\), so \(x = \frac{1}{2}, 5\). The hardest part of this problem
    in my opinion is the factoring.
\end{enumerate}
\end{explanation}

\end{problem}

\begin{problem}

\begin{enumerate}
    \item Write down a quadratic whose zeros are 2 and 5.
    \item Can you find another?
    \item How many such quadratics are there?
\end{enumerate}

\begin{explanation}
\begin{enumerate}
    \item \((x - 2)(x - 5)\) is one example. To check, set equal to zero and solve.
    \item \(-(x - 2)(x - 5)\) is another.
    \item There are infinitely many (just change the constant out front).
\end{enumerate}
\end{explanation}

\end{problem}

\begin{problem}
Alana and Grayson were given the following problem:

\textit{Given a function \(f(x) = x^2 + 7\), find \(f(x - 2)\).}

Alana's work is as follows:
    \begin{align*}
        f(x - 2) & = (x - 2)^2 + 7 \\
        & = x^2 - 2^2 + 7 \\
        & = x^2 - 4 + 7 \\
        & = x^2 + 3
    \end{align*}

Grayson's work is as follows:
    \begin{align*}
        f(x - 2) & = (x - 2)^2 + 7 \\
        & = x^2 - 4x + 4 + 7 \\
        & = x^2 - 4x + 11
    \end{align*}

Whose work is correct? Have a discussion with a partner. Try to ask
each other why.
\begin{explanation}
Grayson's work is correct. When evaluating \((x - 2)^2\), Alana needs to 
remember that squaring is multiplying 2 copies of \(x - 2\), so 
\((x - 2)^2 = (x - 2)(x - 2)\). When multiplying two polynomials, you 
multiply each term in the first by each term in the second and add them (FOIL is a
commonly taught way to remember this in this case).

Thus, \((x - 2)^2 = x^2 - 2x - 2x + 4 = x^2 - 4x + 4\), which is what Grayson wrote.
\end{explanation}
\end{problem}

\end{document}
